\chapter{统计与线性代数}

\section{回归模型}
\noindent\textbf{要点与方法.}\enspace
含截距的回归可先中心化，再用协方差或正交投影求系数。新增解释变量时，把响应与新增变量对旧设计同时残差化，可将多元问题化为一维问题。设计满秩与残差方向非零，分别保证所求系数可识别。

测量误差、样本筛选会改变协方差结构；先写清观测模型和条件总体，再区分线性最优预测与条件均值预测。工具变量须与目标信号相关、与相应误差不相关。投影恒等式属于代数结论，而有限样本的精确 $t$ 检验还需要正确模型下独立、同方差的正态误差。
\par\smallskip



\problem{A01}{测量误差与工具变量}{总体回归；识别}
设 $T,U,V,\varepsilon$ 相互独立、均值均为零且有有限方差，$s=\Var(T)>0$、$u=\Var(U)>0$。实际观测到
\[
 W=T+U,\qquad Z=T+V,\qquad Y=\beta T+\varepsilon.
\]
求 $Y$ 对 $W$ 的总体线性回归斜率。只利用 $(W,Z,Y)$ 的协方差，能否识别 $\beta$？若对同一个 $T$ 取得 $k$ 次具有独立同分布测量误差的读数并取平均，普通回归的偏差如何变化？所有回归含截距。
\solution 独立性使交叉协方差消失，故
\[
 \Var(W)=s+u,\quad \Cov(W,Y)=\beta s,
 \quad \Cov(Z,W)=s,\quad \Cov(Z,Y)=\beta s.
\]
因此普通回归斜率为 $\beta s/(s+u)$，绝对值向零收缩；除 $\beta=0$ 外，它不等于潜在变量的系数。由于 $s>0$，工具变量比值
\[
 \boxed{\beta=\frac{\Cov(Z,Y)}{\Cov(Z,W)}}
\]
有定义且识别目标。这里 $Z$ 与真实解释变量相关，同时与测量误差和响应噪声不相关，两个条件缺一不可。
若 $W_j=T+U_j$，其中各 $U_j$ 与 $(T,\varepsilon)$ 独立且彼此独立、方差为 $u$，则 $\bar W=T+\bar U$，误差方差变为 $u/k$。相应斜率为
\[
 \beta_k=\frac{\beta s}{s+u/k},\qquad
 \beta_k-\beta=-\frac{\beta u}{ks+u}\longrightarrow0.
\]
平均测量消除了误差方差的一部分，却没有把同一个 $T$ 变成 $k$ 个独立的潜在样本。
\insight 区分“回归系数能计算”与“结构系数能识别”；先列可观测协方差，再寻找能消去未知方差的比值。

\problem{A02}{增加变量的拟合增益}{正交投影；部分决定系数}
固定设计 $X\in\R^{n\times k}$ 满列秩且含截距，新增列 $z\notin\operatorname{col}(X)$，$n>k+1$。记仅用 $X$ 时的残差平方和为 $S_0$，使用 $(X,z)$ 时为 $S_1>0$。令 $M=I-X(X^\top X)^{-1}X^\top$、$r=Mz$。求新增系数、$S_0-S_1$，并证明新增变量的部分决定系数与通常 $t$ 统计量的关系。给出精确 $t$ 检验成立的模型假设。
\solution 对固定新增系数 $b$，先最小化原有系数，残差平方和成为 $\|M y-br\|^2$。因 $r\ne0$，配方得到
\[
 \widehat b=\frac{r^\top y}{r^\top r},\qquad
 S_0-S_1=\frac{(r^\top y)^2}{r^\top r}=:D.
\]
这个下降量非负，恰在 $r^\top y=0$ 时为零。它比较的是新增方向与旧残差的对齐程度，而不是 $z$ 与原始响应的相关程度。
设 $\nu=n-k-1$，用 $\widehat\sigma^2=S_1/\nu$ 估计误差方差，则新增系数的标准误和统计量为
\[
 \operatorname{se}(\widehat b)=\sqrt{\frac{S_1}{\nu r^\top r}},
 \qquad t^2=\frac{\nu D}{S_1}.
\]
由于 $S_0=S_1+D$，部分决定系数满足
\[
 \boxed{R^2_{\rm partial}:=\frac{D}{S_0}
       =\frac{t^2}{t^2+\nu}}.
\]
上述代数恒等式不需要随机模型。在正确设定的模型 $y=X\alpha+bz+\varepsilon$ 中，若条件于设计有 $\varepsilon\sim N(0,\sigma^2I)$、$\sigma^2>0$，则在 $b=0$ 下 $t\sim t_\nu$：投影系数的正态分量与正交残差独立，后者的平方和除以 $\sigma^2$ 为 $\chi^2_\nu$。没有正态性时，不能据此宣称有限样本精确 $t$ 分布。
\insight 新增一列只增加一个正交方向；拟合提升、系数检验与部分解释比例因此可以用同一个投影量表示。

\problem{A03}{同号筛选后的预测}{条件分布；预测风险}
$X,Y$ 为独立标准正态变量，只保留事件 $S=\{XY>0\}$。在条件总体中，分别求用 $X$ 预测 $Y$ 的最优含截距线性函数和最优平方损失预测函数，比较二者均方误差。
\solution 标准正态的符号与绝对值独立，且 $\Prob(S)=1/2$。筛选后两个符号相同、正负各半，绝对值仍相互独立且均服从半正态分布。因而条件均值为零、条件方差均为一，且
\[
 \E[XY\mid S]=\E|X|\,\E|Y|=\frac2\pi.
\]
所以最优线性预测为 $\widehat Y_{\rm lin}=2X/\pi$，截距为零，其风险为 $1-4/\pi^2$。
给定 $X=x\ne0$ 与 $S$ 后，只确定了 $Y$ 的符号，故
\[
 m(x)=\E[Y\mid X=x,S]
 =\operatorname{sign}(x)\sqrt{2/\pi}.
\]
这是所有可测预测函数中的平方损失最优解，因为对任意 $g$，条件平方误差可分解为条件方差加 $(m(X)-g(X))^2$。半正态方差为 $1-2/\pi$，故其风险为 $1-2/\pi$，比线性预测低
\[
 \frac2\pi-\frac4{\pi^2}>0.
\]
$x=0$ 是零概率点，预测函数在该点的取值不影响风险。
\insight 筛选可以制造相关性，也可以使条件均值变成非线性函数。先求条件分布，再区分线性最优与总体最优。

\par\Needspace{8\baselineskip}
\subsection*{小节练习}
\exercise{OA02}{残差回归的系数}
设 $Y,X_1,X_2$ 有有限二阶矩，$X_1,X_2$ 的方差均为正，且相关系数 $\rho$ 满足 $|\rho|<1$。所有回归均为含截距的总体最小二乘回归。先将 $Y$ 对 $X_1$ 回归，记残差为 $e$；再将 $e$ 对 $X_2$ 回归，斜率记为 $\beta_1$。在 $Y$ 对 $(X_1,X_2)$ 的联合回归中，$X_2$ 的系数记为 $\beta_2$。求 $\beta_1$ 与 $\beta_2$ 的关系。

\exercise{OA03}{三角区域内的回归}
设 $X,Y$ 相互独立且均服从 $U[0,1]$。在条件 $X+Y>1$ 下，将 $Y$ 对 $X$ 作含截距的总体线性回归。求回归斜率。


\section{参数估计}
\noindent\textbf{要点与方法.}\enspace
估计题先明确比较范围：线性无偏估计、受约束最小二乘与后验估计使用的目标并不相同。协方差矩阵正定时，可用乘子法求最优权重；二次损失配方后，约束的代价表现为特定度量下的距离。相关读数不能按独立观测重复计算信息。

含绝对值的惩罚在零点不可直接求导，应分符号讨论或使用次梯度。共线变量宜先求合计作用，再检查系数分配是否唯一。硬约束与有限软惩罚也须区分。高斯观测可先拆成含信号与纯噪声的方向；本节由不相关推出独立的步骤依赖联合正态性。
\par\smallskip



\problem{A04}{相关观测的均值估计}{广义最小二乘；最优权重}
观测 $Y_i=\mu+\varepsilon_i$，$i=1,\ldots,n$，其中 $n\ge2$、$\E\varepsilon_i=0$，且
\[
 \Cov(\varepsilon_i,\varepsilon_j)=\sigma^2\rho^{|i-j|},
 \qquad \sigma^2>0,\quad0\le\rho<1.
\]
在所有线性无偏估计中求方差最小者及其方差。说明其与算术平均的区别，并求 $\rho\uparrow1$ 时的极限。
\solution 记相关矩阵为 $R$。直接相乘可验证 $R^{-1}=T/(1-\rho^2)$，其中 $T$ 为三对角矩阵：两个端点对角元为 $1$，内部对角元为 $1+\rho^2$，相邻非对角元为 $-\rho$；$n=2$ 时没有内部项。
线性估计 $w^\top Y$ 无偏当且仅当 $\mathbf1^\top w=1$。用乘子最小化 $w^\top Rw$，得 $w=R^{-1}\mathbf1/(\mathbf1^\top R^{-1}\mathbf1)$。端点的行和是 $1/(1+\rho)$，内部行和是 $(1-\rho)/(1+\rho)$。故令 $D=2+(n-2)(1-\rho)$，有
\[
 \widehat\mu=\frac{Y_1+Y_n+(1-\rho)\sum_{i=2}^{n-1}Y_i}{D},
 \qquad
 \boxed{\Var(\widehat\mu)=\frac{\sigma^2(1+\rho)}{D}}.
\]
$R$ 正定保证解唯一。也可将任意其他无偏权重写成 $w+h$、$\mathbf1^\top h=0$；交叉项为零，方差额外增加 $\sigma^2h^\top Rh\ge0$。
$\rho=0$ 时恢复算术平均。$n>2$ 且 $\rho>0$ 时，端点获得更高权重，因为内部观测与两侧均相关；$n=2$ 时仍是等权。固定 $n$ 而 $\rho\uparrow1$，权重趋于两个端点各 $1/2$，最小方差趋于 $\sigma^2$，并不会因读数较多而趋于零。
\insight 重复次数不能替代独立信息量。最优权重由协方差矩阵的逆决定，无须假设误差正态。

\problem{A05}{共线变量的惩罚估计}{Elastic Net；唯一性}
设 $\|x\|^2=1$、$z=x^\top y$，两个解释变量都等于 $x$。对 $\lambda>0$、$\eta\ge0$，求
\[
 \begin{aligned}
 \min_{b_1,b_2}\quad&\frac12\|y-x(b_1+b_2)\|^2\\
 &+\lambda(|b_1|+|b_2|)+\frac\eta2(b_1^2+b_2^2).
 \end{aligned}
\]
的全部最优解。说明为什么极小的正二次惩罚就会改变系数分配的唯一性。
\solution 令 $s=b_1+b_2$。固定 $s$ 时，三角不等式给出 $|b_1|+|b_2|\ge|s|$，平方和满足 $b_1^2+b_2^2\ge s^2/2$。取 $b_1=b_2=s/2$ 同时达到二者下界。因此问题化为
\[
 \min_s\ \frac12(s-z)^2+\lambda|s|+\frac\eta4s^2,
\]
省略项与 $s$ 无关。分别在 $s>0,s<0$ 求导并检查零点次梯度，得
\[
 s^*=\frac{\operatorname{sign}(z)(|z|-\lambda)_+}{1+\eta/2}.
\]
若 $\eta>0$，平方和取等强制 $b_1=b_2=s^*/2$，故系数唯一。若 $\eta=0$，系数和仍唯一；当 $s^*>0$ 时，全部解为 $(t,s^*-t)$、$0\le t\le s^*$；当 $s^*<0$ 时为 $s^*\le t\le0$；当 $s^*=0$ 时，$\lambda>0$ 强制两系数都为零。
因此 Lasso 在非零拟合下可以任意分配两个相同变量的系数；任意 $\eta>0$ 则选择对称分配，同时对总系数施加额外收缩。边界 $|z|=\lambda$ 属于零解情形。
\insight 对共线变量先固定它们的合计作用，再优化内部的系数分配，能将二维问题分成两步一维问题。

\problem{A06}{线性约束下的最小二乘}{约束投影；惩罚极限}
设 $X$ 满列秩，$A=X^\top X$、$\widehat\beta=A^{-1}X^\top y$。给定非零向量 $c$ 与实数 $d$，求约束 $c^\top\beta=d$ 下的最小二乘解，以及相对于无约束解的残差平方和增量。再求 $\lambda\ge0$ 时软惩罚 $\|y-X\beta\|^2+\lambda(c^\top\beta-d)^2$ 的解，并分析 $\lambda\to\infty$。
\solution 正规方程使交叉项消失，故
\[
 \|y-X\beta\|^2=\|y-X\widehat\beta\|^2
 +(\beta-\widehat\beta)^\top A(\beta-\widehat\beta).
\]
记 $q=c^\top A^{-1}c>0$、$e=c^\top\widehat\beta-d$。乘子方程表明校正量沿 $A^{-1}c$ 方向，代入约束得
\[
 \boxed{\widehat\beta_c=\widehat\beta-\frac e q A^{-1}c},
 \qquad \Delta\mathrm{RSS}=\frac{e^2}{q}.
\]
正定二次型保证该可行解唯一。$e=0$ 时约束没有增加代价。
对软惩罚解 $\beta_\lambda$，一阶条件为
\[
 A(\beta_\lambda-\widehat\beta)
 +\lambda c(c^\top\beta_\lambda-d)=0.
\]
左乘 $c^\top A^{-1}$ 可得约束误差为 $e/(1+\lambda q)$，从而
\[
 \beta_\lambda=\widehat\beta-
 \frac{\lambda e}{1+\lambda q}A^{-1}c
 \longrightarrow\widehat\beta_c.
\]
有限惩罚通常不精确满足约束；除 $e=0$ 外，误差仅按 $1/\lambda$ 的量级消失。
\insight 约束投影使用的是 $A$ 定义的距离，而不是普通欧氏距离。约束误差与方向代价分别由 $e$ 和 $q$ 控制。

\problem{P34}{相关传感器的信息合并}{高斯条件分布；充分统计量}
$X\sim N(\mu,s^2)$，$s^2>0$。两个传感器给出 $Z_1=X+\varepsilon_1$、$Z_2=X+\varepsilon_2$。噪声与 $X$ 独立，联合正态、均值为零、方差均为 $\tau^2>0$，相关系数为 $-1<\rho<1$。求 $X$ 的后验分布，证明两读数的差不提供额外信息，并分析 $\rho$ 接近两个端点时的情形。
\solution 令 $M=(Z_1+Z_2)/2$、$D=Z_1-Z_2$。则 $M=X+\bar\varepsilon$，其中
\[
 v:=\Var(\bar\varepsilon)=\frac{\tau^2(1+\rho)}2,
 \quad \Var(D)=2\tau^2(1-\rho).
\]
噪声方差相同，故 $\Cov(\bar\varepsilon,D)=0$；又有 $\Cov(X,D)=0$。联合正态性使 $D$ 与 $(X,M)$ 独立，而 $(M,D)$ 与 $(Z_1,Z_2)$ 一一对应。因此条件分布仅依赖平均读数 $M$。
先验密度乘以 $M$ 的似然，指数中关于 $x$ 的部分为
\[
 -\frac12\left\{\frac{(x-\mu)^2}{s^2}
                  +\frac{(M-x)^2}{v}\right\}.
\]
配方得
\[
 X\mid Z_1,Z_2\sim N(m_*,V_*),\qquad
 V_*=(s^{-2}+v^{-1})^{-1},\quad
 m_*=V_*\left(\frac\mu{s^2}+\frac Mv\right).
\]
$\rho\uparrow1$ 时 $v\to\tau^2$，两读数趋于重复同一个噪声观测，信息量退化为单个传感器。$\rho\downarrow-1$ 时 $v\to0$，平均噪声相消，后验方差趋于零；在端点模型中 $M=X$ 几乎处处成立。对于固定 $s^2,\tau^2$，$V_*$ 随 $\rho$ 严格增加。
\insight 先把多重观测分解成含信号的方向与纯噪声方向，再做条件推断。相关性会决定两个读数究竟互补还是重复。


\par\Needspace{8\baselineskip}
\subsection*{小节练习}
\exercise{OA05}{Lasso 的零解条件}
给定 $X\in\R^{n\times p}$、$y\in\R^n$，其中 $n\ge1$。考虑
\[
 L(\beta)=\frac1{2n}\|y-X\beta\|^2+\lambda\|\beta\|_1,
 \qquad \lambda\ge0.
\]
这里没有另设未受惩罚的截距。求 $\beta=0$ 是最优解的充要条件。


\section{矩阵与二次型}
\noindent\textbf{要点与方法.}\enspace
相关矩阵与 Gram 矩阵都是半正定矩阵。迹、秩及 Rayleigh 商把相关性、向量夹角与特征值问题联系起来；局部元素合法不保证整个矩阵相容。证明界后，还须构造取等对象，并允许奇异的边界情形。

只有矩信息时，可优化平移后的平方矩界，再检查取等分布是否属于所求事件；严格事件可能只有上确界。对充分统计量条件化可使方差不增，完备性可进一步确立无偏类中的最优性。

简单假设间的最强检验由似然比排序，并按原假设校准大小。正态收缩比较的是总平方风险；其维数与协方差条件必须核对，不能据此宣称每个坐标都得到改善。
\par\smallskip



\problem{P07}{缺失相关系数的范围}{相关矩阵补全；构造}
三个随机变量 $X,Y,Z$ 均已标准化。已知 $\Corr(X,Y)=a$、$\Corr(X,Z)=b$，其中 $|a|,|b|<1$。求 $c=\Corr(Y,Z)$ 的全部可行值，并为每个可行值构造联合分布。代入 $a=3/5,b=4/5$，解释区间端点的退化关系。
\solution 令 $U=Y-aX$、$V=Z-bX$。它们与 $X$ 不相关，且
\[
 \Var(U)=1-a^2,\quad\Var(V)=1-b^2,
 \quad\Cov(U,V)=c-ab.
\]
Cauchy--Schwarz 不等式给出必要条件
\[
 \boxed{|c-ab|\le\sqrt{(1-a^2)(1-b^2)}}.
\]
反过来，设 $q=(c-ab)/\sqrt{(1-a^2)(1-b^2)}\in[-1,1]$。取独立标准正态 $G_0,G_1,G_2$，并令
\[
 \begin{aligned}
 X&=G_0,\\
 Y&=aG_0+\sqrt{1-a^2}G_1,\\
 Z&=bG_0+\sqrt{1-b^2}\bigl(qG_1+\sqrt{1-q^2}G_2\bigr).
 \end{aligned}
\]
直接计算即得三个方差均为一且相关系数符合要求，故条件充分。$|q|<1$ 时构造矩阵可逆，相关矩阵正定；$q=\pm1$ 时两个残差线性相关，矩阵奇异。
数值情形给出 $ab=12/25$、根号项也是 $12/25$，故 $0\le c\le24/25$。下端点满足 $V=-3U/4$，上端点满足 $V=3U/4$；这说明达到端点不是相关矩阵失效，而是三个变量之间出现精确线性关系。
\insight 剔除共同变量以后，未知相关性变成两个残差的协方差；必要界与可达构造可以在同一坐标系中完成。


\problem{A25}{给定方差的单侧尾界}{矩约束；取等分布}
设 $\E X=0$、$\Var(X)=\sigma^2>0$。给定 $t>0$，求 $\Prob(X\ge t)$ 的最大可能值及所有取等分布。若改为严格事件 $X>t$，最大值是否仍能达到？
\solution 对任意 $a\ge0$，事件 $X\ge t$ 包含于 $\{(X+a)^2\ge(t+a)^2\}$，故
\[
 \Prob(X\ge t)\le\frac{\E(X+a)^2}{(t+a)^2}
 =\frac{\sigma^2+a^2}{(t+a)^2}.
\]
对 $a$ 求导，导数的符号由 $at-\sigma^2$ 决定，所以最小上界在 $a=\sigma^2/t$ 处取得，得到
\[
 \boxed{\Prob(X\ge t)\le\frac{\sigma^2}{\sigma^2+t^2}}.
\]
取等时，事件外的平方项必须为零、事件内必须恰等于阈值，因此只能在 $t$ 与 $-\sigma^2/t$ 两点取值。均值为零确定 $\Prob(X=t)=\sigma^2/(\sigma^2+t^2)$，另一点取剩余概率；直接核算二阶矩为 $\sigma^2$，故界可达且取等分布唯一。
严格事件的概率仍小于或等于同一界，但若达到，必同时在非严格界取等；唯一取等分布却在 $t$ 没有计入严格尾部，矛盾。因此无最大值。将上述构造中的上点换成 $t+\varepsilon$、下点换成 $-\sigma^2/(t+\varepsilon)$，严格尾概率趋于该界，所以它是上确界。
\insight 对 $X$ 先作平移再使用平方矩，可以保留单侧信息；平移量也是需要优化的变量。

\problem{A29}{单位向量的相干度}{Gram 矩阵；迹不等式}
设 $m>d\ge1$，$u_1,\ldots,u_m$ 为 $\R^d$ 中的单位向量。证明
\[
 \max_{i\ne j}|u_i^\top u_j|^2\ge\frac{m-d}{d(m-1)}.
\]
给出取等所需且充分的两个条件，并在 $d=2,m=3$ 时构造取等向量。
\solution 令 $G$ 为 Gram 矩阵。它半正定、秩至多 $d$，且 $\operatorname{tr}G=m$。把特征值补零至 $d$ 个，Cauchy--Schwarz 给出
\[
 \operatorname{tr}(G^2)\ge\frac{m^2}{d}.
\]
另一方面，若 $M=\max_{i\ne j}|u_i^\top u_j|^2$，则
\[
 \operatorname{tr}(G^2)=m+\sum_{i\ne j}|u_i^\top u_j|^2
 \le m+m(m-1)M.
\]
联立即得所求。两步同时取等，当且仅当 $G$ 恰有 $d$ 个非零特征值且都为 $m/d$，并且所有非对角内积的绝对值相同。第一条件等价于
\[
 \sum_{i=1}^m u_i u_i^\top=\frac md I_d.
\]
取二维中方向角为 $0,\pi/3,2\pi/3$ 的三个单位向量，非对角内积的绝对值都为 $1/2$，平方为 $1/4=(3-2)/(2\cdot2)$，达到界。
\insight 平方内积总和由 Gram 矩阵的谱控制，再由平均值逼出至少一对较大的内积。这里控制绝对值，不能只研究内积的正负。

\problem{A30}{秩一更新的特征值}{极小极大原理；谱扰动}
实对称矩阵 $A$ 的特征值为 $\lambda_1\le\cdots\le\lambda_n$。令 $B=A+tuu^\top$，其中 $t\ge0$、$u\in\R^n$，记其特征值为 $\mu_1\le\cdots\le\mu_n$。证明当 $i<n$ 时
\[
 \lambda_i\le\mu_i\le\lambda_{i+1},
\]
并证明 $\mu_n\ge\lambda_n$、每个特征值增量至多 $t\|u\|^2$、所有增量之和恰为 $t\|u\|^2$。
\solution 每个向量上的 Rayleigh 商都增加 $t(u^\top x)^2/\|x\|^2\ge0$，由极小极大公式得 $\mu_i\ge\lambda_i$。同时增加量不超过 $t\|u\|^2$，再次使用同一公式即得 $\mu_i\le\lambda_i+t\|u\|^2$。
为证更强的交错上界，取 $A$ 前 $i+1$ 个特征向量张成的空间 $E$。当 $u\ne0$ 时，$E\cap u^\perp$ 的维数至少为 $i$；在这个交空间上更新项为零、$A$ 的 Rayleigh 商至多 $\lambda_{i+1}$。选其中一个 $i$ 维子空间代入 $B$ 的极小极大公式，得到 $\mu_i\le\lambda_{i+1}$。$u=0$ 或 $t=0$ 时结论直接成立。
最后取迹，得到
\[
 \sum_{i=1}^n(\mu_i-\lambda_i)
 =\operatorname{tr}(tuu^\top)=t\|u\|^2.
\]
因此更新的总谱位移有固定预算，而各个有序特征值都不会下降。$n=1$ 时交错部分为空，其余结论仍成立。
\insight 秩一扰动在一个余维为一的子空间上完全消失；这比单纯估计扰动范数提供更细的特征值位置控制。

\problem{CT01}{二阶概率的无偏估计}{充分统计量；Rao--Blackwell}
$X_1,\ldots,X_n$ 独立服从 $\operatorname{Bernoulli}(p)$，$n\ge2$、$0<p<1$。从无偏估计 $X_1X_2$ 出发，构造 $p^2$ 的方差更小或相等的无偏估计，求其精确方差，并证明它在所有无偏估计中方差最小。
\solution 令 $S=\sum_iX_i$。给定 $S=s$，成功位置在所有大小为 $s$ 的子集中等概率分布，故
\[
 U:=\E[X_1X_2\mid S]=\frac{S(S-1)}{n(n-1)}.
\]
条件分布不依赖 $p$，所以 $S$ 是充分统计量。全方差公式给出 $\Var(X_1X_2)=\Var(U)+\E\Var(X_1X_2\mid S)$，保证无偏性不变且方差不增加。
还可将 $U$ 写成所有 $\binom n2$ 个乘积 $X_iX_j$ 的平均。单项方差为 $p^2(1-p^2)$，共享一个指标的两项协方差为 $p^3(1-p)$，不相交两项独立。共享指标的无序项对有 $n\binom{n-1}{2}$ 个，因此
\[
 \Var(U)=\frac{2p^2(1-p^2)+4(n-2)p^3(1-p)}{n(n-1)}.
\]
为证明最小方差，先证完备性：若 $\E_p h(S)=0$ 对所有 $p\in(0,1)$ 成立，除以 $(1-p)^n$ 后得到关于 $p/(1-p)$ 的恒零多项式，所以每个系数 $\binom ns h(s)$ 都为零。
对任意另一无偏估计 $W$，其条件均值 $\E[W\mid S]$ 是不含未知参数的函数，且与 $U$ 的期望相同；完备性迫使二者相等。因此再次用全方差公式得 $\Var(W)\ge\Var(U)$。无有限方差的无偏估计也不可能改善该界。$n=2$ 时 $U=X_1X_2$，方差公式仍成立。
\insight 条件化去掉与参数无关的随机性；完备性进一步保证没有其他无偏的充分统计量函数可供改进。

\problem{CT02}{指数样本的最强检验}{似然比；精确功效}
$X_1,\ldots,X_n$ 独立服从率参数为 $\theta$ 的指数分布。检验 $H_0:\theta=\theta_0$ 对 $H_1:\theta=\theta_1>\theta_0$。在第一类错误概率至多 $\alpha\in(0,1)$ 的检验中，求最强检验及其功效。给出 $n=2,\theta_0=1,\theta_1=2,\alpha=1-2/e$ 时的明确临界值与功效。
\solution 令 $T=\sum_iX_i$。备择与原假设的联合密度比为
\[
 L=\left(\frac{\theta_1}{\theta_0}\right)^n
      e^{-(\theta_1-\theta_0)T},
\]
随 $T$ 严格递减，故候选检验为 $\phi_*=\ind_{\{T<c\}}$，其中 $c>0$ 唯一满足 $\Prob_{\theta_0}(T<c)=\alpha$。$T$ 连续，无须在边界随机化。
设相应似然比阈值为 $k$。任何其他检验 $0\le\phi\le1$ 都有 $(\phi_* -\phi)(L-k)\ge0$。若其大小不超过 $\alpha$，则
\[
 \E_{\theta_1}(\phi_* -\phi)
 \ge k\E_{\theta_0}(\phi_* -\phi)\ge0.
\]
所以候选检验的功效最大，这也给出了似然比原则的证明。
指数变量之和的密度为 $\theta^nt^{n-1}e^{-\theta t}/(n-1)!$；可由卷积归纳得到。逐次分部积分给出
\[
 \Prob_\theta(T<c)=1-e^{-\theta c}
       \sum_{j=0}^{n-1}\frac{(\theta c)^j}{j!}.
\]
用 $\theta_0$ 解临界值、再用 $\theta_1$ 代入即为功效。指定数值下 $c=1$，功效为 $1-3e^{-2}$。率参数越大，等待总时间倾向越短，因此拒绝域位于左尾。
\insight 先用似然比确定应选择哪一侧，再用原假设分布校准大小；检验方向不能只凭参数名称猜测。

\problem{CT03}{正态均值的收缩风险}{Stein 恒等式；平方损失}
观测 $X\sim N_d(\theta,I_d)$，$d\ge3$，未知 $\theta\in\R^d$。在估计族
\[
 \delta_a(X)=\left(1-\frac a{\|X\|^2}\right)X
\]
中求使平方损失风险最小的常数 $a$，并证明它对每个 $\theta$ 都优于直接使用 $X$。在 $X=0$ 处任意定义估计，不影响风险。若用同一估计预测独立的 $Y\sim N_d(\theta,I_d)$，其预测风险如何变化？
\solution 写 $\delta_a=X+g(X)$，其中 $g(x)=-ax/\|x\|^2$。对高斯密度分部积分可得
\[
 \E[(X_i-\theta_i)g_i(X)]=\E[\partial_i g_i(X)].
\]
本题的零点奇性不会破坏等式：先在 $\|x\|>\varepsilon$ 积分，球面边界项的量级为 $\varepsilon^{d-2}\to0$；导数的局部量级为 $\|x\|^{-2}$，在 $d>2$ 时可积，无穷远边界由高斯衰减消失。
直接求导得到
\[
 \|g(x)\|^2=\frac{a^2}{\|x\|^2},\qquad
 \sum_i\partial_i g_i(x)=-\frac{a(d-2)}{\|x\|^2}.
\]
展开风险并应用上述恒等式，得
\[
 \begin{aligned}
 R_\theta(a)&=\E_\theta\|\delta_a(X)-\theta\|^2\\
 &=d+\bigl(a^2-2a(d-2)\bigr)
                 \E_\theta\frac1{\|X\|^2}.
 \end{aligned}
\]
最后的期望有限且严格为正，所以对任意 $\theta$，最优常数同为 $a=d-2$。风险为
\[
 d-(d-2)^2\E_\theta\frac1{\|X\|^2}<d,
\]
而直接估计 $X$ 的风险为 $d$。这比较的是全部坐标的平方误差总和，并不声称每个坐标分别改善。
独立性及 $\E(Y-\theta)=0$ 使预测交叉项消失，所以 $\E\|Y-\delta_a(X)\|^2=d+R_\theta(a)$；最优收缩常数不变，最优预测风险严格小于直接预测 $X$ 所得的 $2d$。
\insight 允许适当偏差可以降低总平方风险。高斯分部积分将难处理的交叉误差改写为收缩函数的散度。


\par\Needspace{8\baselineskip}
\subsection*{小节练习}
\exercise{OA30}{主子矩阵的特征值}
设 $A$ 为 $n\times n$ 实对称矩阵，$n\ge2$。删去 $A$ 的第 $k$ 行和第 $k$ 列得到矩阵 $B$。分别将特征值按升序排列为
\[
 \lambda_1\le\cdots\le\lambda_n,
 \qquad \mu_1\le\cdots\le\mu_{n-1}.
\]
证明两组特征值之间的交错关系。


\section{练习题}
\noindent\textbf{要点与方法.}\enspace
这些练习要求把模型条件转成可核对的等式与不等式。观测矩题先检查同一潜在模型能否同时满足全部协方差关系；极值概率题则把“证明上界”与“达到上界”分开处理，严格事件的边界质量尤其重要。

约束估计先列可行域，再检查乘子符号、互补松弛与唯一性。预算把坐标联系起来，不能只做各自的阈值运算。得到候选解后，可用直接代入、单变量消元或目标值增量作独立核对，最后说明结论适用的参数范围。
\par\smallskip



\pagebreak
\exercise{AG1-1}{观测矩的相容性}
沿用 A01 的独立零均值模型，观测矩为
\[
 \begin{aligned}
 \Var(W)&=5,\quad\Var(Z)=4,\quad\Cov(W,Z)=3,\\
 \Cov(W,Y)&=\Cov(Z,Y)=6.
 \end{aligned}
\]
识别 $\beta$、$T,U,V$ 的方差，并给出 $Y$ 对 $W$ 的回归斜率。若只把 $\Cov(Z,Y)$ 改为 $5$，其余观测矩不变，模型还能成立吗？

\exercise{AG1-2}{严格尾事件的最优界}
在所有 $\E X=0$、$\Var(X)=1$ 的分布中，分别求 $\Prob(X\ge2)$ 的最大值与 $\Prob(X>2)$ 的上确界。写出前者的取等分布，以及逼近后者上确界的一列分布。

\exercise{AG1-3}{预算约束的稀疏估计}
中心化设计满足 $X^\top X/n=I_3$、$X^\top y/n=(3,1,-1)^\top$。求
\[
 \min_{b_j\ge0,\ \sum_jb_j\le2}
 \frac{\|y-Xb\|^2}{2n}+\frac12\sum_{j=1}^3b_j.
\]
给出预算约束的乘子，并解释为什么仅对各坐标独立截断还不够。

\exercise{AG1-4}{系数和约束的代价}
一个满列秩最小二乘问题满足 $X^\top X=\operatorname{diag}(1,4)$，无约束解为 $(2,-1)^\top$。要求两个系数之和为零，求受约束解以及残差平方和的增加量。用直接单变量最小化再核对一次。
